\chapter{General introduction}


\section{Why investigate subject focus?} \label{sub:1.1}
\largerpage
Language variation can be defined as the availability of two or more forms for the same grammatical, morphological, or phonological phenomenon within the same variety \parencite{labov1972sociolinguistic}. The choice between these forms, however, is not a matter of pure optionality. Rather, variation is determined by specific choices, often reflecting particular communicative needs. Our utterances are indeed organized in such a way that the speaker adds or omits information based on his/her own intention and on the (presumed) hearer’s knowledge of the current state of affairs. In other words, communication is  organized and structured according to discourse-domain principles, which are reflected by different linguistic means and vary from language to language \parencite{krifka2008basic, komagata2003contextual, wardbirner2001discourse}. The term \textit{Information Packaging} \parencite{chafe1976givenness} or \textit{Information Structure} \parencite{halliday_1967} indicates how information conveyed by linguistic means fits into the (hearer’s mental model of the) context or discourse.

Speaker and hearer share a discursive \textit{Common Ground} (henceforth CG) \parencite{stalnaker1974pragmatic, karttunen1974presupposition, lewis1979scorekeeping, simons2025availability}, where information is constantly managed through communicative strategies, such as questions and answers. \textcite{roberts1996information} believes that CG-management is driven by hierarchies of open questions in need to be settled. Within the universe of discourse, the new, relevant piece of information that resolves an open question is generally referred to as \textit{focus}. Focus is a phenomenon that originates in discourse structure, semantics, and pragmatics (\cite{krifka2008basic, rooth1985association, rooth1992theory, schwarzschild1999givenness, von1990current}). 

Within the alternative-semantics framework \parencite{rooth1985association, rooth1992theory}, focus is defined as that part of the sentence selected from a series of contextually relevant alternatives that fulfill the wh-variable contained in an explicit or implicit question. 
When such a new, relevant piece of information corresponds to the grammatical or semantic subject of the sentence, the literature generally employs the term \textit{subject focus}. An example is provided in (\ref{ex:intro:1}):

\ea \label{ex:intro:1}
\ea Who opened the door?
\ex \ [Mary]\textsubscript{\textsc{foc}} \normalsize [opened it]\textsubscript{\textsc{background.}}\normalsize
\z \z

In (\ref{ex:intro:1}), the subject \textit{Mary} satisfies the wh-variable \textit{which} is contained in the question and is therefore analyzed as the focus. The rest of the clause (`opened it') represents information already shared in the CG by speaker and hearer, and is generally referred to as the \textit{background} proposition \parencite{von1982structured, krifka1991compositional, krifka2015structured}.

 Focused subjects are by no means an understudied topic. Over the years, the phenomenon has attracted the attention of several scholars. Such interest stems most probably from the frequently observed correlation between this informa\-tion-structural function and a plethora of syntactic, prosodic and morphological phenomena.\footnote{Note that some approaches consider the mapping between focus and syntax a correlation. Focus exists first as a discourse concept, and syntax/prosody reflects it. For example, the fact that focus tends to appear in certain syntactic positions (like sentence-final or cleft constructions) is treated as a statistical or functional tendency, not a strict rule. Conversely, according to other approaches, focus is not only correlated to, but rather defined by, formal correlates (e.g., \cite{chomsky1970some}) - that is, it is only by having a formal marker (such as stress placement, syntactic movement, clefting) that something counts as focused. I would like to thank Andreas Dufter for this suggestion.} This correlation has been reported for many closely related and unrelated languages, where focus on the subject usually triggers non-canonical phenomena, while focus on a non-subject argument is generally said to be expressed through \say{unmarked,} canonical strategies.
 
This difference in encoding focus has been defined as the \say{argument asymmetry} in focus realization \parencite{skopeteas2010focus}, and it can be summarized as follows: If, for some languages, focus on subjects obligatorily induces a non-canonical marking strategy, focus on non-subjects does so only optionally. Consequently, 
``if a non-canonical structure occurs with focus on non-subjects, it is expected to occur with focus on subjects too” \parencite[170--171]{skopeteas2010focus}. Evidence for the existence of this phenomenon has been provided for Hausa \parencite{hartmann2007place}, West Chadic languages \parencite{zimmermann2011grammatical}, several Kwa and Gur languages \parencite{fiedler2005out}, and Northern Sotho \parencite{zerbian2007investigating}.

Romance languages, such as French and Spanish, are no exception to this tendency. \textcite{lambrecht1996information}, in his seminal work on the mapping of information-structural categories onto syntactic forms, claimed that both French and Spanish speakers make use of a combination of complex syntactic and prosodic operations to prevent the hearer from interpreting the subject of the sentence as the topic. More specifically, in French, a subject interrogative (e.g., `Who opened the door?') is commonly answered using a rather heavy syntactic strategy, namely a cleft sentence, while in Spanish, the same context would trigger a reordering of the constituents of the clause, with the subject appearing in a postverbal position:

 \ea \label{ex:intro:2}  {French}
\ea \textit{Qui a ouvert la porte?}\\
 ‘Who opened the door?'
\ex \textit{C'est}  {[\textit{Marie}}]\textsubscript{\textsc{foc}} \normalsize  \textit{qui l'a ouverte}.\\
`It is Marie who opened it.'

\z \z 

 \ea \label{ex:intro:3} {Spanish}
\ea \textit{¿Quién abrió la puerta?}\\
 ‘Who opened the door?'
\ex \textit{La abrió} {[\textit{Maria}]\textsubscript{\textsc{foc.}}}\normalsize\\
`Maria opened it.'

\z \z 

Following this tradition, a considerable number of articles investigating narrowly focused subjects in French and Spanish have been produced by both formal and functional schools (for Spanish: \cite{zubizarreta1998prosody, costa2001postverbal, samek2001crosslinguistic, gutierrez2013structural, gabriel2007fokus, gabriel2010focus, heidinger2013information, uth2014spanish, cabrera1999funciones, feldhausen2015oraciones}; French: \cite{prince1978comparison, sornicola1988clefts, dufter2009focus, apotheloz2012pseudo, blanche1990franccais, roubaud2000constructions, fery2001focus, choi2005gens, lahousse2006, Lahousse2011, lahousse2012word, destruel2012french, karssenberg2017oiseaux, Karssenberg+2018}).\footnote{A more comprehensive list of references can be found in Chapter 2 of this book.} Many of these authors have questioned Lambrecht's (\citeyear{lambrecht1996information}) one-to-one correspondence between form and function, arguing that the claim that focus on the subject is encoded in French by clefts and in Spanish by postverbal subjects is probably too simplistic and categorical: Spanish may also use cleft constructions, though less freely than French \parencite{cabrera1999funciones, feldhausen2015oraciones}:

 \ea \label{ex:intro:4}
 \textit{Es} {[\textit{Juan}]}\textsubscript{\textsc{foc}} \normalsize \textit{el que viene}.\\
 `It is Juan who comes.' \hfill \parencite[8]{feldhausen2015oraciones}

 \z 
 

Similarly to Spanish, French may use (a type of) subject-verb inversion for a restricted set of sentences (see \cite{Lahousse2011} for a detailed description):

 \ea \label{ex:intro:5}
 \textit{Paieront une amende} {[\textit{tous les automobilistes en infraction}}]\textsubscript{\textsc{foc.}}\normalsize\\
 `All drivers in breach of the law will pay a fine.'  \\
 \hfill  \parencite[109]{marchello2008evolution}

 \z
 
Such findings have repercussions for a typology of focus realization. Many authors tend to categorize languages based on their primary focus-marking strategies. However, language-internal variation is often neglected, especially in theoretical accounts.
The empirical elicitation of focus, on the other hand, has provided fertile ground for discussion. For instance, a noteworthy finding is that focus on a subject argument in both languages can also be expressed by (prosodically marked) in situ subjects (e.g., \cite{gabriel2007fokus}, \cite{fery2001focus}, \cite{destruel2012french}), as in (\ref{ex:intro:6}) and (\ref{ex:intro:7}):

 \ea \label{ex:intro:6} Spanish
 \ea \textit{¿Quién compró el diario?}\\
 `Who bought the newspaper?'
 \ex  {[\textit{María}}]\textsubscript{\textsc{foc}} \normalsize \textit{compró el diario}.\\
 `María bought the newspaper.' \hfill \parencite[2]{gabriel2007fokus}   
  \z \z

 \ea \label{ex:intro:7} French
 \ea \textit{Qui a escaladé la montagne?}\\
 `Who climbed the mountain?'
 \ex  {[\textit{Arnim}}]\textsubscript{\textsc{foc}}
 \normalsize \textit{a escaladé la montagne}.\footnote{Note that not all scholars accept that preverbal subjects in French can realize subject focus (e.g., \cite{lambrecht1996information}).}\\
 `Arnim climbed the mountain.'  \hfill \parencite[10]{fery2001focus}   
\z \z
 

The fact that subject focus in both languages can be encoded using only pros\-o\-dy shows that Lambrecht's (\citeyear{lambrecht1996information}) categorizations are often too strict and that a more flexible classification is desirable.

Even within the empirical framework, however, the majority of the studies so far produced do not exhaust the formal possibilities offered by the two languages mentioned. For instance, in experimental studies that elicit focused subjects, speakers are usually instructed to avoid answering with a single constituent and to reply rather with sentences containing a verb (e.g., \cite{skopeteas2010focus}).

Yet, in spoken natural language, both French and Spanish speakers have access to a variety of different alternatives, whose alternation is possibly determined by several contextual and linguistic factors. This is the case of fragment answers \parencite{merchant2005fragments} and so-called ``reduced” clefts \parencite{belletti2008answering, doetjes2004cleft, hedberg2000referential, buring1998identity, mikkelsen2012so}, visible in (\ref{ex:intro:8}) and (\ref{ex:intro:9}), respectively:

\ea \label{ex:intro:8}
\ea Who opened the door? 
\ex \ [Mary]\textsubscript{\textsc{foc.}}\normalsize
\z \z

\ea \label{ex:intro:9}
\ea Who opened the door? 
\ex It was [Mary]\textsubscript{\textsc{foc.}}\normalsize
\z \z


Despite receiving less attention, these strategies are nonetheless very frequent in spontaneous speech. Their configuration is somehow less complex and more economical than sentences where both the focus constituent and the background proposition are expressed. In the functionalist tradition \parencite{Givon1985, Givon1995, Haiman1985, macwhinney1989functionalism}, language structure is an iconic reflection of external factors. This means that ``length, complexity, or interrelationship of elements in a linguistic representation parallels the form, length, complexity, or interrelationship of elements in the concept, experience, or communicative strategy that representation encodes” \parencite[756]{newmeyer1992iconicity}. Adopting this functionalist view, it is natural to wonder when speakers would choose more material and when they use less to express the same communicative function.  According to \textcite{Givon1985}, language economy is taken to be a major motivation for iconicity, so any iconic motivation is by definition an economical one.

This book also considers strategies such as those illustrated in (8) and (9). Building on spontaneous speech data, the book provides a comprehensive account of how subject focus is expressed in French and Spanish declarative sentences and investigates the appropriateness conditions for the use of alternative focus-marking devices for the same focus category. Concepts such as language economy are crucial to account for the alternation among different forms that express the same communicative function.

The main goal is, therefore, to investigate whether the occurrence of different marking strategies for focused subjects is motivated by specific linguistic and/or pragmatic factors, by adopting a quantitative variationist perspective \parencite{labov1972sociolinguistic}. Some of these factors have already been discussed by both formal and functional approaches in relation to focus: For instance, it has been claimed that the type of focus, and especially the notion of \textit{contrast}, plays a fundamental role in determining the syntactic form produced. In her seminal work on Spanish focus realization, \textcite{zubizarreta1998prosody} claims that subject foci must obligatorily occur in the postverbal position. However, in the presence of contrast or correction \parencite{van2004incompatibility, steube2001correction, REPP20101333}, the focused subject might also be expressed in situ. Similarly, different cleft formats in French have been said to correlate with different types of focus: \textit{C'est}-clefts, for instance, are generally reported as expressing exhaustive focus (see \cite{kiss1998identificational} for the notion of exhaustivity), while \textit{il y a}-clefts would occur more frequently in contexts where exhaustivity is not conveyed \parencite{de2015status, karssenberg2017oiseaux}. 

It has also been claimed that the choice among different cleft formats is linked to the givenness \parencite{firbas1966non, chafe1970meaning, prince1992zpg} and the weight of the subject constituent. \textcite{lambrecht1996information} advocates a strong correlation between \textit{il y a}-clefts and newly introduced subjects, while \textcite{roubaud2000constructions} has noted the frequent occurrence of heavy clausal subjects with pseudo-clefts\footnote{The term can be attributed originally to \textcite{bach1968english}.}, as in (\ref{ex:intro:10}):

 \ea \label{ex:intro:10}
  \textit{Ce qui m'énerve c'est} {[\textit{quand il crie}}]\textsubscript{\textsc{foc.}}\normalsize\\
 `What bothers me is when he screams.' \hfill \parencite[178]{roubaud2000constructions}    
\z


Additionally, it has been said that focus interacts with the argument structure of the verb in languages like Spanish or Italian \parencite{burzio1986italian, rizzi1982issues, belletti1988case}, and that this interaction has an impact on the position of the subject. \textcite{gutierrez2013structural} attributes this phenomenon to the differences in the thematic properties of the arguments in the sentence, whereas \textcite{gabriel2007fokus} claims that the number of arguments itself is responsible for the presence of more than one focus-marking strategy : If transitive verbs favor in situ focusing of the subject, monoargumentals are said to trigger postverbal positioning (VS order).

 While systematic analyses involving focused subjects and their correlation with verb type (or argument structure) have been carried out for Spanish, the same cannot be said for French. More specifically, focus on the subject in French has not yet been investigated in relation to the argument structure of the verb accompanying the subject. 
Building on previous studies, this project aims to fill this gap by investigating the semantic and contextual factors that lie behind speakers' choice among different strategies, showing that language-internal syntactic variation is not free and unconstrained but rather regulated by specific communicative and linguistic conditions.

A second problem with previous accounts is of a methodological nature. Many functionally oriented authors investigating information structure have pointed out the need for highly comparable material based on spontaneous speech. Information structure, being the way in which communication is ``tailored” \parencite{prince1981towards} and organized between speaker and hearer, is by definition more tangible in dialogical spontaneous situations than in written or formal texts (e.g., narratives, journalistic texts). Thus, a cross-linguistic study of different pragmatic functions would require highly comparable data. 

This book addresses the issue of comparability by building on the analysis of two highly comparable sub-corpora of spontaneous dialogues (\textit{sgs}corpus), collected by \textcite{adli2011gradient} in 2005. Both the French and the Spanish sub-corpora have been elicited using the same game task and exhibit, as a consequence, many similarities in the setting or the content of the dialogues. The investigation of information structure has always presented a methodological challenge. At present, scholars working empirically in this area agree that the method of the question-answer pair is the most reliable way to retrieve the correct information-structural partition of the sentence (e.g., Question Under Discussion approach to focus; \cite{roberts1996information, riester2018annotation, riester2019constructing}). The \textit{sgs}corpus, being mostly dialogical, presents a methodological advantage in this respect.

Finally, for both French and Spanish focused subjects, many authors to date have generalized their claims to the language in its entirety, not taking into account the presence of diatopic variation among regions. In the case of Spanish, at least, information-structural phenomena are particularly sensitive to such variation.\footnote{In this respect, the two languages differ considerably: While for Spanish, diatopic variation has been claimed to affect the way in which focus is realized, in French, the phenomenon seems to be rather homogenous, and no noticeable regional differences have been reported. I thank Andreas Dufter for this suggestion.} One example is the case of Southern Peninsular Spanish (SPS, including Andalusian or Extremaduran), where non-contrastive focus on the subject is said to be expressed more often in a preverbal, canonical position rather than postverbally \parencite{jimenez2015focus}. 

 \ea \label{ex:intro:11}
\ea \textit{El chocolate que tenía escondido ya no está. ¿Quién lo ha encontrado?}\\
 ‘The chocolate I had hidden is not there anymore. Who has found it?'
\ex \ {[\textit{Jimena}}]\textsubscript{\textsc{foc}} \normalsize \textit{lo encontró. Y se lo ha comido entero}.\\
`Jimena found it. And she has eaten it all.' \\
\hfill \parencite[6]{jimenez2015focus}
\z \z 


Considering this, the results for the corpus studies presented in this book are delimited to a very specific variety of both French (Parisian French) and Spanish (Barcelona Spanish), and the claims are restricted to those specific varieties exclusively. Consequently, they may not fully apply in a different geographical area. 

Overall, this book is meant to make a further contribution to the intricate question of focused subjects and their realization in French and Spanish by adopting a comparative variationist approach. It does so by building on spontaneous-speech data, considering every conversational strategy through which this function is expressed, and providing an account for the pragmatic and linguistic motivations behind speakers' choices to opt for one strategy rather than another.


\section{Goals and structure of the book} \label{sub:1.2} 

This book is divided into three major parts: Part I contains the Theoretical Background, where I summarize previous accounts of focused subjects. I start by delimiting the notion of focus and its expression, reviewing at first existing literature on focus as a universal information-structural category
\parencite{zimmermann2011focus}, then narrowing down to subject focus expression. I concentrate particularly on the case of constituent focus realization and on the different focus types that have been identified in the literature (information, contrastive, etc.) (\textsc{\chapref{sec:1}}).

I then present both formal and functional approaches to focused subjects in Spanish and French, first summarizing the studies that have been conducted on the two languages both individually and comparatively (\textsc{\chapref{sec:2}}).

Part II includes three comparative corpus studies carried out using the same corpus (\textit{sgs}corpus, \cite{adli2011gradient}), but tackling different aspects of subject focus realization. As has been often said, both clefting in French and constituent re-ordering in Spanish are but two of several means to focus subject constituents and compete constantly with a series of alternative focus-marking devices. Sentence accent and phrasing, for instance, where prosody functions as the only means for constituent focus in situ, have been claimed to be valid alternatives to word order \parencite{dufter2009focus}. 
The range of possibilities, however, is not limited to these categories. As illustrated above, more economical strategies, such as fragment answers (8) or reduced clefts (9), can also be used to encode the same pragmatic function. The first research goal of the book is thus to clearly identify the \textit{envelope of variation} (\cite{labov1972sociolinguistic}; i.e., the strategies that speakers produce to mark focused subjects in both languages), and discuss them from both a language-internal and comparative perspective. 

In \textsc{\chapref{sec:3}}, I report the results of the subject-focusing strategies extracted from both the Spanish and French sub-corpora. This part has an exploratory purpose and is mostly descriptive. As will be shown, focus-marking strategies for subjects are relatively infrequent (as already argued by others; see, for instance, \cite{brunetti2009italian} for Spanish and Italian), and some structures show very few occurrences. Given that a subject in focus is, by definition, a marked information-structural configuration, the scarcity of such data points is not surprising. Yet, the corpus exhibits variation in the realization of narrowly focused subjects, displaying a gamut of competing means for both Spanish (e.g., postverbal subjects, in situ, clefting, fragments, reduced clefts) and French (e.g., in situ, clefting, fragments, reduced clefts, dislocations). Interestingly, some structures appear in both languages (e.g., clefts, fragments, in situ). However, their frequency varies considerably from one language to the other.

\textsc{\chapref{sec:4}} investigates the alternation between ``full” (e.g., clefts) and ``elliptical” (e.g., fragment) answers\footnote{The terms \textit{full} and \textit{elliptical} refer here to answers where the background proposition is either uttered (e.g., a cleft) or omitted (e.g., a fragment), respectively.} in both languages.
The corpus query has revealed that speakers use more or less economical means to express a focused subject. Some of these structures, such as postverbal subjects, clefts, and in situ focusing, exhibit a clear focus-background partition, where the subject is the focused element of the sentence and the rest of the clause\footnote{Note that cleft sentences are biclausal structures. The background proposition, in this case, is generally contained in the relative clause.} is the background (information already shared and known by speaker and hearer). Other strategies only display the focused part (the subject), while the background proposition has been elided. This is the case of fragment answers (8) and reduced clefts (9). 

The systematic study of competing variants requires not only the identification of these forms but also the individual contexts in which differences between them may arise. To tackle this question, each strategy is investigated in its context (that is in relation to the preceding explicit or implicit question it constitutes an answer to). The main intuition is that, in a dialogical context, the activation of the background proposition plays an important role when speakers choose between full forms, such as clefts or postverbal subjects, and elliptical forms, such as fragments or reduced clefts. In other words, when the background material is highly activated because the question immediately precedes the answer, speakers opt for a more economical form(a fragment answer). Conversely, when the background material is not fully activated because of intervening utterances or when the question is implicit, a cooperative speaker will use a more explicit form (e.g., a cleft) to facilitate comprehension. The idea that shorter forms, such as reduced clefts, are felicitous with immediately preceding questions has already been postulated in the literature \parencite{belletti2008answering, katz2000categories}. However, it has never been empirically investigated, nor operationalized. 

In \textsc{\chapref{sec:4}}, I propose an operationalization of the distance between question and answer, adopting the Question Under Discussion (QUD) approach to focus \parencite{roberts1996information, buring2003d, beaver2009sense, riester2018annotation} and distinguishing among three degrees of immediacy of the QUD (immediate, distant, implicit). The results show that immediate QUDs favor the presence of elliptical answers, while distant or implicit QUDs are more frequently answered with full forms. Although there is no one-to-one correspondence between level of activation of the backgrounded proposition and strategy used, the results reported in \textsc{\chapref{sec:4}} suggest that the choice of the sentential form is, to some extent, influenced by principles such as language economy, or the \textit{Principle of Communicative Efficiency} \parencite[30]{Levshina_2022}.

\textsc{\chapref{sec:5}} narrows down the investigation to the syntactic phenomena that have mostly been discussed in the literature: Spanish subject position (SV/VS) and different cleft formats. The analysis reported in this chapter investigates the interacting semantic and discourse factors that have been claimed to influence the occurrence of these phenomena. As mentioned in \sectref{sub:1.1}, it has already been claimed that the way in which focus is realized is dependent on other factors, such as the length of the constituent, type of focus conveyed, or verb type (e.g., \cite{gutierrez2013structural}, \cite{gabriel2007fokus}). If a considerable number of publications have found an effect of argument structure in focus expression, the same reasoning has never been systematically applied to French, and the role of argument structure in interaction with focus remains unclear for this language. Given that French clefts have been  considered a discourse-pragmatic equivalent of Spanish postverbal subjects \parencite{lambrecht1996information, lambrecht2010constraints, belletti2008answering, lahousse2012word}, the lack of literature investigating this correlation is somewhat surprising.

In \textsc{\chapref{sec:5}}, I show how specific factors influence speakers' preferences of different syntactic structures, albeit to varying degrees. Because of typological differences between the two languages, these factors affect subject focus realization at different levels: While verb type and givenness of the subject result in word-order alternations (SV, VS) in Spanish, the same combination of factors result in different cleft formats in French, a language with a relatively rigid order of constituents.

Finally, Part III presents a comparative analysis of the results for the two languages \textsc{\chapref{sec:6}}. The findings are interpreted from a variationist perspective \parencite{labov1972sociolinguistic, labov1994principles, labov_2006}. Within this framework, a \textit{variable} is understood as a general or abstract feature that is subject to variation, whereas a \textit{variant} is the specific realization of the variable in speech \parencite{meyerhoff2018introducing}. The results of this study show that different focus-marking strategies within the same language can be regarded as optional variants of a single variable and that their distribution in natural speech can be explained by specific conditioning factors.

The main goals of the study are thus to investigate whether there are differences and similarities in the way two closely related languages encode subject focus and to understand the motivations behind speakers' communicative choices. The research questions of the book can be summarized as follows: 

\begin{itemize}
\item \textbf{Research question 1}: What answering strategies are used to mark subject focus in French and Spanish? How do the two languages differ in this respect?
\item \textbf{Research question 2}: What determines speaker choice between full and elliptical forms when expressing focused subjects in the two languages?
\item \textbf{Research question 3}: How do factors such as focus type, givenness, and argument structure contribute to syntactic variation in the expression of subject focus in the two languages?
\end{itemize}

The outcome of this study makes several innovative contributions: i) it offers an empirical contribution to the investigation of subject focus realization from both a language-internal and comparative perspective, ii) it provides an account for the alternation between full and elliptical focus-marking strategies in terms of distance between question and answer, and iii) it investigates factors that govern the syntactic realization of focus, both individually and in interaction.

It is important to keep in mind that the findings of this book are highly register-dependent. Being based on dialogical data, the same results may not be obtained if applied to, say, written and/or formal texts (e.g., narratives, journalistic texts). This is not surprising if we consider that the pragmatic function of some constructions, such as cleft sentences, varies considerably across different text types. Therefore, the claims issued from the results of the present analysis are not intended as a general theory on on cleft or postverbal-subject usage in any possible contexts. Rather, the investigation is limited to the usage of these constructions in spontaneous dialogues, which, however, are ``one of the most frequent registers experienced by humans'' \parencite[836]{du1987discourse}.\footnote{I would like to thank Aria Adli for this suggestion.}
